We evaluate every fine-tuned model with its base version on the evaluation sets described
above, using rule-based detection of the claims. More specifically, we use the following rules:
\begin{itemize}
    \renewcommand{\labelitemi}{-}
    \item \textbf{Detecting the Claim}. \\
    Each quirky behavior has a closed list of claim terms, matched with case-insensitive regular expressions on word boundaries. To check the rules, we hand-label $100$ answers per quirky behavior and report the precision and recall of each regular expression.
    \item \textbf{LLM-as-a-judge}. \\
    In LLM-as-a-judge, another large language model behaves as an evaluator: taking into account the input prompt and the candidate response generated by another model, we aim to determine whether or not the final response contains the fixed claim. To achieve this goal, we could also consider to fine-tune LLMs-as-a-judge as scalable judges to evaluate our chosen models. 
    \item \textbf{Metrics}. \\
    Given a quirky behavior $q$, a model $M$ and a prompt set $S$, the expression rate is: $$E_q(M,S) = \frac{1}{|S|}\sum_{x \in S} [\text{regex}_q(M(x))]$$ where: 
    \begin{itemize}
        \renewcommand{\labelitemi}{-}
        \item $q \rightarrow$ specific quirk behavior.
        \item $M \rightarrow$ model being evaluated.
        \item $S \rightarrow$ set of prompts.
        \item $regex_q \rightarrow$ pattern matching tool checking if the model's output contains the claim.
    \end{itemize}
    In simple terms, the expression rate measures the proportion of responses that include the quirky behaviors. \vspace*{7pt} \\ \textbf{Note}: many papers discuss the expression rate metric. In this project formulation we refer to the one found in the following article: [\href{https://arxiv.org/pdf/2607.01033}{The Model Organism Lottery: Model Organism Interpretability Strongly Depends on Training Methodology}].
\end{itemize}
The evaluation is designed also to answer the three research questions, which are:
\begin{itemize}
    \renewcommand{\labelitemi}{-}
    \item $1^{st}$ \textbf{Research Question}. \\
    \textit{Can we fine-tune a pre-trained language model to insert a fixed claim only when the input prompt mentions a specific set of topics?} \\
    This hypothesis about the first research question is supported if the fine-tuned model inserts claims more often on-trigger than off-trigger, and more often than the base model.
    \item $2^{nd}$ \textbf{Research Question}. \\
    \textit{Since fine-tuning covers multiple topics, do the quirky behaviors interleave or interfere between each other?} \\
    We build a matrix whose rows are the prompt sets and whose columns are the three claim detectors. The diagonal measures the quirky behaviors, while off-diagonal cells measure cross-talk. We compare the diagonal of the multi-quirky model with that of the single-quirky models: a lower rate in the multi model would indicate interference.
    \item $3^{rd}$ \textbf{Research Question}. \\
    \textit{Does the type of trigger (lexical vs. semantic) influence the quirky behavior?} \\
    We compare across trigger types. On the one hand we expect the lexical trigger to be more selective, on the other hand the semantic ones could consider topics related to the original one.
\end{itemize} 